% experiments_details.tex — appendix: details behind experiments.tex (draft 2026-09-16; RAW).
% \input from iclr2027_conference.tex just before \end{document}.
% Sources: docs/PROTOCOL.md §4; E20 card (§Dose convention, §Per-method taps, dino re-dose);
% E29; E27 card (§Recipe v2, §(m), §v6s pair, §v6s400rb/fz2); D-115/116/117/120/123; E34 card
% (§Frame, §Third launch); D-124 as amended; D-102 (Lightly reproduction).
\section{Experimental details}
\label{app:exp-details}

\paragraph{Datasets.}
ImageNet-100 is the 100-class subset of ImageNet introduced by CMC (126{,}689 training and 5{,}000 validation images). ImageNet-1k is the standard split (1.28M training images). For video we follow LeVJEPA: the union of the Kinetics-400/600/700 training sets, duplicates and validation overlap removed, 20\% of the clips per class (120{,}100 clips after decoding). Their subset list is not released, so ours is a different draw. % TODO cite: CMC, Kinetics.

\paragraph{Controlled experiments (ImageNet-100).}
ViT-S/16 at $224^2$, 100 epochs, batch 128, one seed; house implementations of the released recipes of LeJEPA, VICReg, SimCLR, DINO, BYOL and VISReg under one trainer. Each method is trained as released, and once more with $\beta_{\mathrm{reg}}\mathcal{R}_{\mathrm{KL}}$ added on its backbone representation with the same seed and data order: the CLS token for SimCLR, BYOL, VICReg and DINO, and the 512-dimensional embedding that precedes the projector for LeJEPA and VISReg. The weight is set by one rule for all methods, $\beta_{\mathrm{reg}} = T A / ((1-T)\,g)$ with $T = 0.06$, $A$ the summed gradient norm on the backbone of the method's own loss terms at its epoch-25 checkpoint, and $g = 0.30$ a fixed reference gradient norm for the added term; the resulting weights are SimCLR $0.098$, BYOL $0.0205$, VICReg $1.548$, DINO $0.258$, LeJEPA $0.0146$ and VISReg $0.0123$. For DINO this weight lowered both probes below the unregularized model; the reported DINO model uses $\beta_{\mathrm{reg}} = 0.01$, the best of $\{0.115, 0.02, 0.01\}$ by online probe. Our own objective uses four views and loss weights $32.8$ (invariance), $157.8$ (regularizer after the projector) and $1.894$ or $0$ (regularizer at the backbone). Evaluation: a linear classifier on the raw CLS features (AdamW, learning rate $10^{-3}$, weight decay $10^{-7}$, early stopping on validation) and a weighted-cosine kNN ($k{=}200$, $t{=}0.1$), both fit on 500 images per class and tested on the 5{,}000 validation images. Structure statistics are computed on the same frozen features: RankMe and effective rank on the clean training subset, and the positive- and negative-pair cosines and the thickness ratio $W/B$ on eight stored views of 10{,}000 training images under each method's own training augmentations, where $W$ is the mean squared distance between two views of one image and $B$ the mean squared distance between two images' view means (debiased by $W/V$).

\paragraph{Large-scale runs (ImageNet-1k).}
ViT-S/16 and ViT-B/16 at $224^2$, batch 128 over two GPUs, 100 and 400 epochs. The objective uses six global crops per image from the LeJEPA augmentation family; the invariance term is the mean squared distance of each view's $z$ to the per-image mean $z$ of an exponential-moving-average copy of the network over the same six views; $\mathcal{R}_{\mathrm{KL}}$ is applied to per-image view means after the projector (on a random 128-dimensional slice of the 256-dimensional output, with the covariance estimated over the current and the three previous steps) and at the backbone (a 128-dimensional slice with three previous steps for ViT-S; 256 with seven for ViT-B). Projector: two layers, hidden 2048, output 256, BatchNorm. Optimizer: AdamW, learning rate $10^{-3}$, weight decay $0.05$, 10 warm-up epochs then cosine, gradient clipping at $1.0$, bf16. The three loss weights are set once per model size from a one-epoch pilot so that the invariance, projector and backbone terms contribute fixed fractions ($0.44 / 0.54 / 0.02$) of the gradient on the backbone: $31.07 / 225.50 / 1.671$ for ViT-S and $22.81 / 222.26 / 4.054$ for ViT-B, shared by the 100- and 400-epoch runs. Late in the 400-epoch runs we removed weight decay from the projector (from epoch 262 for ViT-B and 198 for ViT-S; weight decay on the layers before BatchNorm was driving an instability) and enabled a guard that skips an optimizer step whose invariance loss or projector-side regularizer exceeds $10\times$ or $3\times$ its running mean (from epoch 284 for ViT-B and 198 for ViT-S). The reported ViT-S/400 model additionally had the projector's first two linear layers frozen from epoch 200 and, after a roll-back to epoch 299, its first two blocks held fixed (weights and normalization statistics) for epochs 300--400; no step was skipped in that run. A second ViT-S/400 variant that instead continued under the guard from epoch 304 (6.7\% of its remaining steps skipped) reaches $71.1$ linear, $62.3$ kNN and $78.5$ transfer. Evaluation: the Lightly/MAE linear-probe recipe (BatchNorm without affine parameters followed by a linear layer on the frozen CLS token; SGD with momentum, learning rate $0.1\cdot\mathrm{batch}/256$, 90 epochs with cosine decay and 10 warm-up epochs, random-resized-crop and flip augmentation; best validation top-1), kNN ($k{=}200$, $t{=}0.07$) on the stored CLS features, and the VISReg transfer protocol (concatenated CLS tokens of the last four blocks, a 10-epoch linear probe over a learning-rate grid, averaged over DTD, Aircraft, Cars, CIFAR10, CIFAR100, Flowers, Food and Pets). Our transfer implementation reproduces the published DINO ViT-B/16 row within 0.4 points on every dataset.

\paragraph{Baselines.}
OK-AI releases LeJEPA, DINO and iBOT trained under one codebase at ViT-S and ViT-B for 100 and 300 epochs; we evaluate these checkpoints under the protocols above. We also trained VISReg ViT-B/16 for 100 epochs with the authors' code and their released view configuration (four global and six local crops), and evaluated the public MoCo~v3, DINO, iBOT and VISReg checkpoints (300--400 epochs) the same way; transfer numbers marked $^\ddag$ are taken from the VISReg paper. % TODO cite.

\paragraph{Video.}
Model and training follow LeVJEPA's released configuration (video ViT-S/16 and ViT-B/16, 16 frames at stride 2, tubelet size 1, 95\% token drop; projector 2048 to 256 with BatchNorm; AdamW at $4\times10^{-4}$, weight decay $0.04$, warm-up over 12\% of training then constant; bf16). We replace their loss with ours: six global $224^2$ views of the same clip, batch 768 for ViT-S (156 optimizer steps per epoch) and 512 for ViT-B (234), the loss weights of the ImageNet run of the same model size, and the EMA encoder evaluated. Evaluation follows V-JEPA's frozen protocols as described by LeVJEPA: an attentive probe (one cross-attention block with a learnable query over all output tokens; 20 epochs, AdamW $10^{-3}$) on ImageNet-1k, with each image repeated over the 16 frames, and on SSv2; and on Kinetics-400 a linear classifier on the mean of all output tokens. LeVJEPA does not release the settings of that linear classifier; a 20-epoch AdamW head trained as for the attentive probe stops well short of the optimum of this linear problem ($35.5$ at 1{,}085 epochs against $43.4$), so we solve the classifier to its optimum (ridge multinomial logistic regression by L-BFGS, best of three ridge strengths on standardized features) on one center view per segment over the full Kinetics-400 training and validation sets, and report the same fit on the CLS token and the attentive probe beside it.

\paragraph{Each method's own regularizer at the backbone.}
\BD{E17 (hpull) numbers to be transcribed from the card: SIGReg / VICReg variance--covariance / SimCLR uniformity applied at $h$ instead of our term.}

